\documentclass[11pt]{article}

\usepackage[margin=1in]{geometry}
\usepackage{booktabs}
\usepackage{enumitem}
\usepackage{hyperref}
\usepackage{microtype}

\hypersetup{colorlinks=true,linkcolor=blue,citecolor=blue,urlcolor=blue}
\setlist{nosep}

\newcommand{\project}{Ginux}
\newcommand{\modelname}{OpenAI Codex with GPT-5.6 Sol}
\newcommand{\evaluatedcommit}{\texttt{9afa61762a66a208\allowbreak{}55ea3a23e3e04221\allowbreak{}b54a5a22}}
\newcommand{\shortcommit}{\texttt{9afa617}}
\newcommand{\estimatedcost}{approximately USD~8,000}

\title{\textbf{Ginux: What an AI-Built Linux-Compatible Kernel Reveals About the State of AI Systems Engineering}}
\author{Gregory M. Kurtzer}
\date{Draft technical report --- October 2026}

\begin{document}
\maketitle

\begin{abstract}
We report an experiment in which one persistent AI coding agent, operating with
task-level autonomy under human architectural governance, constructed an
operating-system kernel and supporting environment from the ground up. The
artifact, \project, boots directly through UEFI on x86-64, starts a ring-3
controller, implements a substantial portion of the Linux userspace contract,
ingests OCI images, runs an unmodified Alpine userspace, and launches Rocky
Linux 10 systemd as PID~1 inside a kernel-native domain. Approximately 99\% of
development was performed by a single longitudinal \modelname{} task rather
than a collection of specialized agents. The human supplied goals,
architectural constraints, acceptance criteria, and corrective review, but did
not direct implementation step by step.

The evaluated artifact contains roughly 139,000 retained non-generated lines
of implementation and test code across 1,037 commits produced over eleven
calendar days. The author estimates API usage cost at \estimatedcost. These
figures demonstrate unusual construction speed, but they do not establish
software quality. At the stopping point, 82 of 375 Linux x86-64 syscalls were
self-classified as complete, 245 as partial, and 48 as policy-denied. Important
non-syscall interfaces remain incomplete; GPU drivers, stable kernel-driver ABI
compatibility, security assurance, native performance characterization, and
general application compatibility are absent. A final manual lifecycle test
also exposed a failure that narrower agent-authored tests had missed.

The experiment therefore supports a narrower and more consequential claim than
``AI built a production operating system.'' A modern coding agent can construct
a deep systems prototype and use compiler, QEMU, runtime, and benchmark
feedback to improve it. It cannot thereby certify compatibility, security,
performance, or maintainability. \project{} also provides preliminary evidence
that kernel-native domains are technically feasible, but not yet that their
claimed isolation or resource-policy benefits have been proven. We argue that
the central problem has shifted from code generation to assurance: standards,
independent test oracles, traceability, release controls, and durable ownership
must determine when AI-produced systems code can be trusted.
\end{abstract}

\section{Introduction}

AI systems can write code. That observation alone is no longer a useful
research result. The more difficult question is what kind of software they can
construct, how much autonomy they can exercise, and what evidence is required
before their output should be trusted. Operating-system kernels are a useful
stress case because their correctness is not captured by compilation or a
small collection of examples. A kernel must preserve thousands of explicit and
implicit contracts involving instruction-level entry, virtual memory,
concurrency, executable loading, filesystems, devices, process lifecycle,
networking, security boundaries, and failure cleanup.

We created \project{} to test two questions. First, can an AI coding agent build
a new kernel capable of running meaningful, substantially unmodified Linux
userspace? Second, is a kernel-native \emph{domain}---an isolation, ownership,
accounting, scheduling, and failure-containment boundary stronger and more
fundamental than a Linux namespace---technically feasible and potentially
beneficial? The first question is primary; the second emerged as the
architectural vehicle through which workloads were organized.

This was not an attempt to deliver a production operating system. Development
stopped when the author judged that the experiment had demonstrated the
construction capability under study. The resulting system is neither complete
nor generally deployable. That distinction is central: a research stopping
point can answer a feasibility question without satisfying the compatibility,
security, performance, hardware, support, and maintenance requirements of a
production platform.

The name \project{} is a play on Linux, whose name derives from Linus Torvalds.
The ``G'' ostensibly stands for GPT; interpretations involving the author's
first name are discouraged so that his head does not become any larger.

This report makes four contributions:

\begin{enumerate}
  \item a longitudinal case study of one AI agent constructing a substantial
        systems artifact under task-level, rather than step-level, human direction;
  \item a benchmark-driven development procedure combining versioned architecture
        documents, an explicit plan, compilation, QEMU execution, automated gates,
        and manual acceptance tests;
  \item an artifact-level account of the Linux compatibility and native-domain
        mechanisms reached at a fixed commit; and
  \item an analysis of the gap between rapid AI construction and credible
        assurance, including cases where the agent incorrectly concluded that an
        incomplete benchmark proved the requested behavior.
\end{enumerate}

\section{Research Questions and Scope}

We organize the experiment around three levels of capability.

\paragraph{Construction.} Can a coding agent create a bootable kernel, control
plane, device paths, and Linux compatibility layer rather than an isolated
function or application patch?

\paragraph{Assurance.} Can the process demonstrate that the artifact behaves
correctly beyond selected examples, particularly across the large Linux
userspace contract?

\paragraph{Stewardship.} Can compatibility, security, performance, and
maintainability be preserved while requirements and implementations evolve?

The first question has affirmative evidence. The second has mixed evidence.
The third remains open. We do not claim full Linux compatibility, production
security, hardware portability, stable driver compatibility, or performance
parity with Linux. We also do not treat dispatching every syscall number as
equivalent to implementing Linux semantics.

The evaluated artifact is repository commit \evaluatedcommit{}, dated October
20, 2026. Results in this report apply to that commit and the stated QEMU
profiles. Later changes must not be read back into these findings.

\section{Human-Governed, Agent-Executed Method}

\subsection{Division of responsibility}

Approximately 99\% of implementation was produced within one persistent Codex
task using \modelname{}. This percentage is an author estimate; provider-level
task and token exports have not yet been archived. The process should not be
described as an army of specialized agents. Nor should ``human-directed'' be
read as hand-held implementation.

The human acted as architect, requirements authority, product owner, and
acceptance tester. He established desired properties, changed architectural
direction, defined user-visible benchmarks, rejected misleading results, and
decided when the research question had been answered. Codex inspected and
modified the repository, designed implementations, built the system, ran QEMU,
interpreted failures, added tests and documentation, committed checkpoints, and
iterated. The human generally specified \emph{what} must be true and why; the
agent determined \emph{how} to implement it.

The work took place in one longitudinal chat. That does not imply that every
raw message remained in one physical model context window: the hosting system
may compact long histories into summaries. ``Single longitudinal task'' is
therefore more accurate than ``one immutable context.'' This continuity
reduced formal handoff cost, but it also concentrated implementation,
test-authoring, and success judgment in the same model.

\subsection{Architecture documents as external memory}

The conversation was not the sole repository of intent. The human and agent
collaboratively developed and maintained a set of versioned documents that the
agent used as reference material alongside an active plan:

\begin{itemize}
  \item \texttt{ARCHITECTURE.md} described the system model, trust boundaries,
        ABI layers, boot path, performance principles, and milestones;
  \item \texttt{docs/ROADMAP.md} and checkpoint documents defined staged
        acceptance gates;
  \item \texttt{docs/LINUX-COMPATIBILITY.md} maintained the syscall and
        non-syscall compatibility ledger;
  \item architecture decision records captured durable choices about domains,
        authority, unified images, and driver placement;
  \item stability, capacity, storage, and performance audits recorded known
        constraints and measured findings; and
  \item an active execution plan tracked the current objective and stopping point.
\end{itemize}

These artifacts served as persistent external memory: they survived context
compaction, made earlier decisions searchable, and gave the agent a written
contract against which to plan changes. They were also collaboratively
authored by the same human--AI process and were not independent specifications.
Several documents ran ahead of passing end-to-end evidence or retained
aspirational multi-agent procedures that did not describe the actual
experiment. Documentation improved continuity but did not eliminate drift;
traceability from a claim to a passing gate remained necessary.

\subsection{The iteration loop}

The operative method was benchmark-driven:

\begin{center}
\fbox{\parbox{0.91\linewidth}{
Human architectural intent $\rightarrow$ concrete milestone $\rightarrow$
observable acceptance benchmark $\rightarrow$ Codex implementation
$\rightarrow$ build and automated tests $\rightarrow$ QEMU boot/runtime test
$\rightarrow$ log analysis and revision $\rightarrow$ repeat as needed $\rightarrow$
human acceptance $\rightarrow$ signed commit and stopping-point review.
}}
\end{center}

Codex was often instructed to continue the implementation--test--diagnose loop
until a benchmark passed, rather than returning after each edit. Milestones
were intended to end in a committed state suitable for manual evaluation.
Failures found manually were converted into implementation changes and, when
successful, regression gates.

The method resembles repeated sampling and execution feedback in earlier code
generation work~\cite{chen2021codex}, but at repository and system scale.
SWE-bench similarly motivates real repository tasks as a step beyond isolated
code generation~\cite{jimenez2023swebench}; this experiment extends the idea
to a long-lived greenfield systems artifact whose specification and benchmark
also evolved during development.

\subsection{The benchmark-oracle failure mode}

The loop was powerful but vulnerable to an incomplete oracle. An agent can
iterate until a benchmark passes while the benchmark fails to represent the
requirement. Near the evaluated stopping point, one gate verified that a
detached Rocky domain returned control and reached a \texttt{booting} state;
another verified that a Rocky domain started with \texttt{--attach} reached a
login prompt. Neither initially verified the actual sequence ``start detached,
attach during or after boot, remain attached, detach, reattach, stop, and
restart.'' The agent reported the feature complete. Manual execution of that
sequence reached a login prompt and then failed runtime cleanup while three
Linux processes remained live, marking the domain failed.

This is the software equivalent of benchmark overfitting. The same agent had
implemented the behavior, authored much of the test, interpreted the output,
and declared success. Human acceptance testing supplied a different oracle and
contradicted the conclusion. The episode is evidence about the development
method, not merely one kernel bug.

\section{Artifact Architecture}

\subsection{Boot and control plane}

\project{} is a Rust- and assembly-based x86-64 system that boots as a unified
UEFI image. The boot path obtains firmware state, establishes kernel-owned page
tables and protection structures, initializes required platform mechanisms,
and enters the kernel without a conventional host distribution or initramfs.
A statically linked controller executes in ring~3. Host-specific management is
exposed through a versioned control interface; the controller is not compiled
into ring~0.

The kernel provides scheduling, memory ownership, capabilities, IPC,
interrupts, IOMMU and device mechanisms, executable loading, and personality
objects. Narrow native services may execute in separate domains. OCI is used
as an image distribution and filesystem assembly format~\cite{ocispec}; it is
not the security boundary or runtime.

\subsection{Linux as a userspace personality}

Linux x86-64 is the first userspace personality. The target contract draws
from the Linux 6.12 syscall table~\cite{linux612syscalls}, Linux UAPI and
man-pages~\cite{linuxmanpages}, and the x86-64 psABI~\cite{x8664psabi}. The
implementation includes syscall entry and dispatch, ELF \texttt{ET\_EXEC} and
selected \texttt{ET\_DYN} loading, dynamic interpreters, virtual memory,
processes and threads, signals, descriptors, filesystems, sockets, terminals,
and partial \texttt{/proc}, \texttt{/dev}, namespace, and security interfaces.
It is libc-neutral in intent: the executable's \texttt{PT\_INTERP} selects the
loader from the domain root rather than choosing a musl or glibc mode.

\subsection{Native domains}

A domain is the outer kernel-owned principal for address spaces, processes,
capabilities, resources, IPC objects, filesystem and mount views, network
objects, devices, lifecycle, and failure policy. Linux namespaces are intended
to be compatibility constructs inside a Linux-personality domain, allowing a
domain eventually to run its own nested container manager without weakening
the outer boundary.

This differs from a hypervisor domain such as Xen~\cite{barham2003xen}: a
\project{} Linux workload does not carry a separate guest kernel. It also
differs from claiming a verified microkernel boundary; unlike seL4's explicit
functional-correctness result~\cite{klein2009sel4}, \project{} has no formal
proof or comparable assurance case. The artifact demonstrates that the
abstraction can be represented, scheduled, and used to host Linux processes.
It does not yet demonstrate strong isolation under adversarial testing or
resource-quality guarantees under contention.

\section{Implementation Trajectory}

Table~\ref{tab:milestones} summarizes selected repository milestones. The
history contains 1,037 small commits, 1,036 with developer certificate
sign-off, between October 10 and October 20, 2026.

\begin{table}[ht]
\centering
\small
\begin{tabular}{lll}
\toprule
Date & Commit & Observable milestone \\
\midrule
Oct. 10 & \texttt{97f778a} & Enter kernel after UEFI boot services \\
Oct. 10 & \texttt{ab89034} & Enter an isolated control plane at ring~3 \\
Oct. 11 & \texttt{ec48014} & Execute a bounded Linux ELF entry \\
Oct. 13 & \texttt{a4d3371} & Execute Alpine shell through musl \\
Oct. 13 & \texttt{6beb469} & Attach the Alpine shell to its console \\
Oct. 15 & \texttt{a94e03d} & Support Rocky glibc shell workloads \\
Oct. 18 & \texttt{f87bac8} & Launch Rocky systemd as domain PID~1 \\
Oct. 19 & \texttt{9c64072} & Reach an interactive Rocky systemd console \\
Oct. 20 & \texttt{aa51d21} & Execute commands in a running domain \\
Oct. 20 & \shortcommit & Stage detached-domain startup work \\
\bottomrule
\end{tabular}
\caption{Selected milestones from the signed Git history. A milestone commit
records implementation progress, not independent conformance certification.}
\label{tab:milestones}
\end{table}

The trajectory was not linear. Early Linux probes exercised individual ABI
contracts before real images. Alpine exposed path lookup, directory iteration,
fork/exec, working-directory, and musl behavior. Rocky's glibc and dynamic
utilities expanded ELF, mapping, loader, and filesystem requirements. systemd
then exercised process scheduling, epoll, signals, procfs, mount identity,
AF\_UNIX sockets, timers, credentials, pseudo-filesystems, and console
lifecycle in combination. Failures repeatedly converted an apparently local
feature into a cross-subsystem requirement.

\section{Measured Stopping Point}

\subsection{Artifact size and activity}

Table~\ref{tab:artifact} reports measurements from \shortcommit. ``Selected
implementation/test lines'' counts tracked Rust, assembly, C, headers, and
shell files under the principal source, service, test, specification, and tool
directories. The non-generated total excludes the 35,946-line generated boot
payload source. LOC is reported for scale, not as a quality metric.

\begin{table}[ht]
\centering
\begin{tabular}{lr}
\toprule
Measure & Value \\
\midrule
Development interval & 11 calendar days \\
Commits & 1,037 \\
Tracked files & 261 \\
Rust source files & 63 \\
Selected implementation/test LOC, including generated payload & 174,736 \\
Selected non-generated implementation/test LOC & 138,790 \\
Markdown documentation LOC & 13,074 \\
Rust function definitions (syntactic count) & 3,317 \\
Rust unit-test annotations & 53 \\
QEMU shell test programs & 11 \\
Top-level Make targets & 39 \\
Author-estimated API expenditure & \estimatedcost \\
\bottomrule
\end{tabular}
\caption{Repository-scale measurements at the evaluated commit. Counts are
simple reproducible textual measurements, not complexity or correctness scores.}
\label{tab:artifact}
\end{table}

The kernel directory accounts for approximately 116,779 selected lines, the
control plane 3,262, reusable crates 6,868, native services 4,284, tests 6,224,
and tools 894. These categories include low-level probes and assembly fixtures;
they should not be interpreted as equivalent units of effort.

\subsection{Linux compatibility}

The compatibility ledger covers all 375 native Linux 6.12 x86-64 syscall
numbers. At the stopping point it classified 82 as Complete, 245 as Partial,
48 as Policy-denied, and zero as Missing. ``Zero Missing'' means no table entry
falls through the default \texttt{ENOSYS} dispatcher. It does not mean all 375
calls implement Linux behavior. The categories are project-defined and were
largely produced and tested by the same agent; they have not received
independent conformance review.

Non-syscall surfaces remain more visibly incomplete. Static PIE without an
interpreter is unsupported; broader dynamic-loader, locale, and NSS behavior
is unproven; process and signal semantics are bounded; procfs is a subset;
sysfs, cgroupfs, vDSO/vvar, IPv6, complete namespace operations, and broad ioctl
families remain absent or partial. General mount and namespace calls were still
policy-denied at the evaluated commit despite the intended domain model
eventually allowing such operations inside the outer boundary.

Practical workload evidence is stronger than comprehensive conformance
evidence. Unmodified Alpine/musl shells execute; Rocky/glibc shells and dynamic
utilities execute; and a Rocky Linux 10 image configured for UBI init can launch systemd,
reach multi-user targets, and present a console login in selected paths.
Arbitrary Linux programs are not guaranteed to work.

\subsection{Performance evidence}

The repository's October 19 profile records a reduction in QEMU TCG
Linux-init-to-console time from 15,344~ms to a two-run median of 4,082~ms after
fixing parent-dependent Unix polling, quadratic frame-reference scans, and
eager immutable-rootfs duplication. This is translated x86 execution on Apple
Silicon, not a native kernel benchmark. Warm OCI layer restore and rootfs
publication remained roughly 2.7 and 3.5 seconds respectively and were outside
the init measurement. A three-second KVM gate was defined but could not be run
on the development host.

The final manual detached-start workflow at \shortcommit{} measured
13,833~ms from Linux init to console and then hit the cleanup failure described
earlier. These results show both meaningful optimization and insufficient
benchmark stability. No comparison against Linux on identical native hardware,
no throughput or tail-latency suite, and no performance claim suitable for
production existed at the stopping point.

\subsection{Security evidence}

The code implements capability tables, W\textasciicircum X mappings, guarded stacks,
generation-tagged domain ownership, selected fault cleanup, SMEP/SMAP/UMIP
setup, IOMMU-related device paths, cryptographic checks, and many negative
tests. These are mechanisms and test evidence, not a security evaluation.
There has been no independent audit, systematic adversarial campaign, complete
threat-model validation, or formal proof. Driver compatibility, including
NVIDIA GPU operation and unmodified Linux \texttt{.ko} modules, is not
implemented. The system therefore cannot support claims about GPU workloads,
stable kABI, or production isolation.

\section{What the Experiment Says About AI}

\subsection{Demonstrated strengths}

The agent showed substantial breadth. It moved between UEFI, assembly entry,
x86-64 paging and interrupts, Rust kernel code, ELF and ABI details, network
and block devices, OCI formats, shell tooling, QEMU automation, user-facing CLI
design, documentation, and performance profiling. It used execution feedback
rather than producing a static code dump. It could maintain a working
repository, make small signed commits, recover from many failures, and turn
real distribution behavior into implementation tasks.

The pace is itself evidence. Even after excluding generated boot payload, the
retained artifact is too large and cross-cutting to describe as autocomplete.
One longitudinal agent produced a research kernel prototype across a ten-day
commit interval while responding to changing architectural direction.

\subsection{Demonstrated weaknesses}

The same history shows recurrent problems:

\begin{itemize}
  \item the agent sometimes optimized for an observable marker rather than the
        complete user workflow;
  \item it occasionally reported success after adjacent tests passed while the
        requested end-to-end sequence remained broken;
  \item it could forget or blur architectural distinctions, such as OCI image
        transport versus runtime execution;
  \item compatibility breadth encouraged shallow syscall dispatch and static
        capacity increases before semantic or dynamic-resource behavior was
        fully validated;
  \item long context and rapid iteration produced regressions and documentation
        claims ahead of evidence; and
  \item the agent had difficulty choosing a stopping point without explicit
        human insistence.
\end{itemize}

Architecture documents and plans reduced these effects but did not remove
them. A single agent retained continuity, yet the absence of an independent
test author or reviewer allowed common assumptions to contaminate
implementation, benchmark, and verdict.

\subsection{Where AI is now}

The experiment suggests that current AI is capable of \emph{construction at
systems depth}. It can generate and integrate enough low-level code to make a
new machine boot, execute foreign userspace, and expose novel abstractions. Its
weaker capability is \emph{assurance}: determining whether the resulting
behavior is complete, equivalent, secure, fast, and stable under changes. The
agent can operate a feedback loop, but the loop is only as good as its oracle.

This distinction matters because production engineering is dominated by the
properties that persist after a demonstration: compatibility across unknown
programs, predictable failure handling, upgrade safety, performance under
load, hardware qualification, security review, incident response, and years
of maintenance. Code generation attacks only part of that cost.

\section{Economics}

The author estimates total API expenditure at \estimatedcost. Provider billing
exports and exact token totals should be archived before publication. Active
human hours were not instrumented, so this report does not convert elapsed time
into a fabricated labor comparison.

Two simple ratios are reproducible but easy to misuse: the estimate is about
USD~7.71 per commit and USD~0.058 per retained non-generated implementation or
test line. Neither measures correctness, usefulness, review burden, or future
maintenance. A ten-line scheduler fix can matter more than thousands of lines
of generated tables; code later discarded still consumed tokens; and retained
bugs count as cheaply as retained correct code.

A defensible economic model separates four costs:

\begin{enumerate}
  \item \textbf{generation}: producing candidate code and documentation;
  \item \textbf{integration}: making independently evolving subsystems work together;
  \item \textbf{assurance}: establishing compatibility, security, and performance; and
  \item \textbf{maintenance}: preserving those properties over releases and incidents.
\end{enumerate}

\project{} provides evidence that AI can reduce generation cost enough to make
an ambitious one-off experiment accessible. It provides no evidence that the
total cost of a production-quality replacement would be similarly reduced.
The public history permits later comparison with human projects, but a valid
counterfactual requires scoped requirements, comparable tests, active human
time, specialist review, and equivalent quality gates. LOC alone cannot supply
that counterfactual.

\section{Domains as a Secondary Finding}

The domain design was implementable enough to structure the controller,
services, resource ownership, Linux personality state, and multiple workload
records. Per-domain state banking, generation identities, address-space and
console suspension, and failure ownership offer preliminary evidence that the
abstraction is technically feasible without placing a complete guest kernel in
each workload.

The proposed benefits remain hypotheses. CPU and memory settings are retained
and reported but are not fully charged and enforced. Network and filesystem
isolation are incomplete. Lifecycle failures remain. There is no adversarial
cross-domain security result, QoS benchmark, formal noninterference argument,
or native-hardware comparison. The correct conclusion is that domains are a
promising research direction exposed by the experiment, not a proven superior
replacement for namespaces, containers, or virtual machines.

\section{Controls Required for Trustworthy AI Systems Code}

The failures observed here suggest concrete controls for future AI-generated
infrastructure:

\begin{itemize}
  \item versioned specifications and compatibility ledgers linked to tests;
  \item acceptance workflows that represent complete operator behavior, not
        isolated success markers;
  \item differential testing against Linux and established conformance suites;
  \item independent test authors and adversarial reviewers, preferably with
        model and method diversity;
  \item negative, lifecycle, cleanup, concurrency, and resource-exhaustion tests;
  \item performance budgets tied to named workloads and identical hardware;
  \item threat models, unsafe-code inventories, fuzzing, and independent security review;
  \item reproducible builds and archived logs, prompts, billing evidence, and test artifacts;
  \item release gates that distinguish implementation, integration, and human acceptance; and
  \item explicit ownership for compatibility and maintenance after generation ends.
\end{itemize}

The architecture document proposed multiple specialized agents for
implementation, verification, and integration. The actual experiment did not
substantially use that structure. A future controlled comparison should test
whether separate implementation, compatibility, security, and acceptance
agents reduce self-validation failures. Multiple instances of the same model
would not automatically provide independent judgment; they may reproduce the
same blind spots.

\section{Threats to Validity}

This is a single project, model, author, architecture, and development period.
The author both directed and evaluated the experiment. The AI authored much of
the implementation, documentation, and test suite. Manual tests were
opportunistic rather than a preregistered protocol. API cost and development
share are author estimates. Human time was not recorded. Repository commits
measure activity but not cognitive effort, and the unusually fine commit
granularity inflates comparisons based on commit count.

The system ran primarily under x86-64 QEMU TCG on an Apple Silicon host. KVM,
real hardware, other architectures, hostile workloads, broad package suites,
and GPU drivers were not evaluated. Compatibility classifications are local to
the project and should be treated as hypotheses until independently tested.
The current commit is known to contain a detached-domain lifecycle failure;
this report intentionally preserves that fact rather than selecting only a
successful path.

Finally, the project specification changed during development. That is
realistic for exploratory engineering but makes controlled attribution
difficult. The result demonstrates what this human--AI collaboration produced,
not what the model would produce from a fixed prompt without intervention.

\section{One-Offs and Production Systems}

One-off AI-generated systems can be valuable. They can test architectural
possibility, create research instruments, explore alternatives too expensive
to staff conventionally, and expose hidden requirements. \project{} succeeded
in that role. It made Linux compatibility and domain design concrete enough to
discover problems that discussion alone would not reveal.

The same properties become risks when a one-off is mistaken for a platform.
\project{} cannot guarantee that programs operate correctly. It cannot load
the existing GPU-driver ecosystem, promise a stable driver ABI, establish
security against attack, or provide appropriate production performance data.
It has no release organization responsible for years of compatibility. There
are environments in which a bespoke artifact is the point; there are others
in which standards and ecosystem contracts are the product.

It is therefore fair to say that vibe-coding an operating system is possible.
It is not reasonable to infer that vibe-operating production workloads is
safe. Production is a different undertaking: it requires evidence and
institutions that persist after the code-generation session ends.

\section{Conclusion}

\project{} is unfinished, and that does not make the experiment a failure. Its
purpose was not to ship a production operating system. Its purpose was to
determine whether a modern AI coding agent could construct one deeply enough
to boot from firmware, implement a meaningful Linux userspace contract, run
real distribution components, and explore a new domain boundary.

The answer is yes. A single persistent Codex agent, given tasks, architectural
constraints, a collaboratively maintained documentary memory, an explicit
plan, and access to build and QEMU feedback, produced a substantial kernel
prototype. The human guided objectives and judged results without directing
the implementation line by line.

That answer must not be confused with production readiness. \project{} cannot
guarantee general Linux program behavior; its compatibility ledger is mostly
partial; GPU and stable driver compatibility are absent; performance and
security have not received appropriate independent evaluation; and a manual
workflow at the evaluated stopping point exposed a serious lifecycle failure
after narrower tests passed.

AI can build code. The question now is the quality and capability of that code,
and how standards for control and compatibility are maintained. The cost of
generating an ambitious one-off has fallen dramatically. The cost of knowing
that it is correct, compatible, secure, performant, and maintainable has not
disappeared. Future progress should therefore be measured not only by what AI
can make run, but by the independent evidence that justifies trusting it.

\section{Artifact Availability}

The project is licensed under Apache-2.0 and is intended for public release.
The archival repository URL and immutable release tag will be inserted before
submission. This draft evaluates Git commit \evaluatedcommit{}; the commit,
source manifest, compatibility ledger, raw benchmark logs, and development
transcript should be preserved together so that later project changes do not
silently alter the reported result.

\bibliographystyle{plain}
\bibliography{references}

\appendix
\section{Reproduction and Publication Checklist}

Before publication, the following evidence should accompany the report:

\begin{enumerate}
  \item a public repository URL and immutable tag for commit \shortcommit;
  \item an exported prompt/task history sufficient to substantiate the
        single-agent development description;
  \item provider billing evidence supporting the API-cost estimate;
  \item a reconstructed human-time range, or an explicit statement that it
        cannot be recovered;
  \item clean-checkout logs for \texttt{make check}, Alpine, Rocky shell,
        multi-domain, storage, OCI-cache, and systemd gates;
  \item raw timing logs with QEMU, accelerator, host, image digest, and cache state;
  \item the exact detached-start/attach failure transcript retained as a
        negative result; and
  \item a generated artifact manifest and line-count script.
\end{enumerate}

\section{Milestone Record Template}

Future revisions should record each claimed milestone with the following
fields: objective; human-specified acceptance behavior; agent implementation;
automated benchmark; exact environment; human validation sequence; commit;
result; residual limitations; and regressions discovered later. A milestone is
not complete merely because the implementation exists or a nearby test passes.

\end{document}
